\documentclass[11pt,a4paper]{article}

% ---- 한글 + 수식 ----
\usepackage{kotex}
\usepackage{amsmath,amssymb,amsfonts,bm}
\usepackage{graphicx}
\usepackage{geometry}
\geometry{margin=25mm}
\usepackage[hidelinks]{hyperref}
\usepackage{enumitem}
\usepackage{tcolorbox}
\usepackage{xcolor}

% 각주를 페이지마다 1번부터 다시 세지 않고 길게 달기 좋게
\renewcommand{\thefootnote}{\arabic{footnote}}

% "한 줄 설명" 박스
\newtcolorbox{readbox}[1][]{
  colback=blue!4, colframe=blue!40!black, boxrule=0.4pt,
  left=6pt,right=6pt,top=4pt,bottom=4pt, #1}

\newtcolorbox{notebox}[1][]{
  colback=orange!6, colframe=orange!55!black, boxrule=0.4pt,
  left=6pt,right=6pt,top=4pt,bottom=4pt, #1}

\title{\textbf{지지흐름 텐서장(SFTF) 쉽게 읽기}\\[4pt]
\large 3D 프린팅에서 ``어느 방향으로 출력할까''를 빠르게 고르는 방법\\[2pt]
\normalsize ── 고등학생을 위한 수식 한 줄씩 풀어 읽기 ──}
\author{}
\date{}

\begin{document}
\maketitle

\begin{readbox}
\textbf{이 글의 목표.} 원래 논문(\emph{Support Flow Tensor Field}, 줄여서
\textbf{SFTF})에 나오는 수식들이 대학원생도 한 번에 이해하기 어렵다는 이야기를
들었습니다. 그래서 이 글에서는 \emph{같은 내용}을, 고등학교 1\,$\sim$\,2학년이
아는 수학(벡터, 내적, 평균, 그래프)만으로 읽을 수 있도록 \textbf{수식 한 줄마다
바로 밑에 우리말 설명}을 붙였습니다. 처음 보는 용어는 \textbf{각주}\footnote{각주(脚註)란
이렇게 페이지 아래쪽에 작은 글씨로 다는 보충 설명입니다. 본문 흐름을 끊지 않으면서
어려운 단어를 따로 풀어 주는 장치입니다.}로 최대한 많이 풀었습니다. 수학이 부담되면
파란 박스(``한 줄 요약'')만 따라 읽어도 전체 흐름을 잡을 수 있습니다.
\end{readbox}

\tableofcontents
\bigskip

%======================================================================
\section{먼저 큰 그림: 무슨 문제를 푸는가}
%======================================================================

\subsection{3D 프린팅과 ``서포트''}

FDM\footnote{FDM(Fused Deposition Modeling, 용융 적층 모델링): 가는 플라스틱 실을
녹여 한 층(layer)씩 쌓아 올려 입체를 만드는 가장 흔한 3D 프린팅 방식입니다. 케이크에
생크림을 짜서 한 겹씩 올리는 모습을 떠올리면 됩니다.} 같은 3D 프린터는 재료를
\emph{아래에서 위로 한 층씩} 쌓아 물체를 만듭니다. 그런데 공중에 떠 있거나 심하게
기울어진 부분(\textbf{오버행}\footnote{오버행(overhang, 처마): 아래에 받쳐 주는 것이
없는데 옆이나 위로 튀어나온 부분. 처마나 발코니처럼요. 너무 기울면 갓 쌓은 재료가
굳기 전에 흘러내려 모양이 망가집니다.})은 그냥 두면 흘러내립니다. 그래서 프린터는
임시 기둥인 \textbf{서포트}(support, 받침)를 함께 출력해 받쳐 주고, 출력이 끝나면
이 서포트를 뜯어 버립니다.

\begin{notebox}
\textbf{핵심 갈등.} 서포트는 \emph{버리는 재료}입니다. 재료비\,·\,출력시간\,·\,뒤처리
수고를 모두 늘립니다. 그러니 서포트가 \textbf{적게} 생기도록 출력하고 싶습니다.
\end{notebox}

\subsection{왜 ``방향''이 중요한가}

같은 물체라도 \textbf{어떤 자세(방향)로 눕혀서 출력하느냐}에 따라 필요한 서포트의
양이 크게 달라집니다. 컵을 똑바로 세워 출력하면 손잡이 아래가 떠서 서포트가 많이
필요하지만, 옆으로 눕히면 줄어들 수도 있습니다. 그래서 ``서포트가 가장 적게 드는
출력 방향''을 찾는 것이 중요한 문제입니다.

이 ``방향''을 수학에서는 \textbf{단위벡터}\footnote{벡터(vector): 크기와 방향을
함께 가진 양. 화살표로 그립니다. \emph{단위}벡터는 길이(크기)가 정확히 1인 벡터로,
``방향만'' 나타낼 때 씁니다.} $n$ 으로 적습니다. 3차원 공간에서 길이가 1인 화살표는
공의 표면(반지름 1짜리 구) 위의 한 점에 대응하므로, 가능한 모든 방향의 모임을
$S^2$\footnote{$S^2$: ``2차원 구면(sphere)''이라는 기호. 3차원 공간 안에 있는,
반지름 1인 공의 \emph{껍질}을 뜻합니다. 지구본 표면 위의 점 하나가 ``한 방향''이라고
생각하면 됩니다.}라고 씁니다. 즉 $n \in S^2$ 는 ``$n$ 은 어떤 한 방향''이라는 뜻입니다.

\begin{notebox}
\textbf{무식한 방법(전수조사)의 문제.} ``그럼 모든 방향을 1$^\circ$씩 다 돌려보며
서포트 양을 재면 되지 않나?'' 맞지만, 방향의 가짓수가 수만 개라서 모든 방향마다
정밀 계산을 돌리면 \emph{매우 느립니다}. 단순한 모양조차 시간이 오래 걸립니다.
\end{notebox}

\subsection{SFTF의 아이디어 한 문장}

\begin{readbox}
\textbf{한 줄 요약.} SFTF는 정답을 직접 푸는 정밀 계산기가 \emph{아니라},
``이 방향들이 유망해 보인다''를 \textbf{아주 빠르고 값싸게} 골라 주는
\textbf{후보 추천기}입니다. 그다음 추천된 \emph{몇 개 방향 주변만} 정밀 검증하면,
전부 다 검사할 때와 비슷한 좋은 답을, 훨씬 적은 계산으로 얻습니다.
\end{readbox}

비유하면, 넓은 바다에서 보물을 찾을 때 바다 전체를 잠수해 뒤지는 대신(전수조사),
먼저 수면 위에서 ``여기 색이 좀 다른데?'' 하는 곳 20군데를 빠르게 찍고(SFTF),
그 20군데 주변만 잠수해 확인하는(국소 정밀검증) 전략입니다.

%======================================================================
\section{준비물: 고등학교 수학으로 충분합니다}
%======================================================================

본격적인 수식에 들어가기 전에, 앞으로 계속 쓰는 도구 세 가지만 짚겠습니다.

\subsection{면, 면적, 법선벡터}

3D 모델은 보통 수많은 \textbf{삼각형 조각}으로 이루어진 껍질(메쉬\footnote{메쉬(mesh,
그물): 3D 물체의 표면을 잘게 쪼갠 삼각형들의 모임. 삼각형이 많을수록(=면 개수가 많을수록)
표면이 매끄럽고 정밀하지만 계산은 무거워집니다.})입니다. 각 삼각형 면 $i$ 마다 두 가지
숫자를 씁니다.

\begin{itemize}[leftmargin=1.6em]
\item $A_i$ : 그 삼각형의 \textbf{면적}(넓이). 큰 면일수록 더 중요하게 다루려고 씁니다.
\item $m_i$ : 그 삼각형의 \textbf{법선벡터}\footnote{법선벡터(normal vector): 면에
\emph{수직}으로 바깥을 향해 튀어나온 단위벡터. 책상 위에 종이를 놓으면 천장을 향하는
화살표가 그 종이의 법선입니다. ``이 면이 어느 쪽을 바라보는가''를 나타냅니다.}.
길이 1짜리 화살표로, 그 면이 ``어느 쪽을 바라보는지''를 알려 줍니다.
\end{itemize}

\subsection{내적: 두 방향이 ``얼마나 같은 쪽''인지 재는 자}

두 단위벡터 $a$, $b$ 의 \textbf{내적}\footnote{내적(dot product, 점곱): 두 벡터를 곱해
\emph{하나의 숫자}를 얻는 연산. 성분이 $a=(a_1,a_2,a_3)$, $b=(b_1,b_2,b_3)$이면
$a\cdot b = a_1b_1+a_2b_2+a_3b_3$. 단위벡터끼리는 $a\cdot b=\cos\theta$($\theta$=두
벡터가 이루는 각)가 되어, 방향이 얼마나 비슷한지를 $-1\sim 1$ 사이 숫자로 알려 줍니다.}
$a\cdot b$ 은 다음과 같이 읽으면 됩니다.

\begin{equation}
a\cdot b = \cos\theta \quad(\theta = \text{두 화살표 사이의 각도})
\end{equation}
\noindent\textit{읽기:} 두 방향이 \textbf{완전히 같으면} $a\cdot b=1$, \textbf{직각}이면
$0$, \textbf{정반대}면 $-1$ 입니다. 즉 내적은 ``두 화살표가 얼마나 같은 쪽을
가리키나''를 하나의 숫자로 압축한 것입니다. 이 한 줄만 기억하면 뒤의 수식 대부분이
풀립니다.

\subsection{평균을 모으면 ``전체 경향''이 보인다}

수많은 면의 정보를 하나하나 보는 대신, \textbf{전부 더해서(합) 전체 경향}을 보는
일을 자주 합니다. 기호 $\sum$ 은 ``모두 더하라''는 뜻입니다\footnote{$\sum_i x_i$
($i$에 대한 합, 시그마): $x_1+x_2+x_3+\cdots$ 처럼 여러 항을 전부 더하라는 기호.
$\sum_{i=1}^{3} x_i = x_1+x_2+x_3$.}. 예를 들어 $\sum_i A_i$ 는 ``모든 삼각형 면적을
다 더하라'' = 표면 전체 넓이입니다.

%======================================================================
\section{출발점: 대칭 지지 텐서 (그리고 PCA 이야기)}
%======================================================================

SFTF를 만들기 전에, 먼저 더 단순한 ``기준선(baseline)''\footnote{기준선(baseline):
성능을 비교하기 위한 \emph{출발점} 또는 \emph{비교 대상}. ``적어도 이것보단 나아야
한다''는 잣대입니다.} 하나를 봅니다. 면의 법선들이 \emph{전체적으로 어느 쪽으로
쏠려 있는지}를 요약하는 행렬입니다.

\begin{equation}
M=\sum_i A_i\, m_i m_i^{T}
\label{eq:M}
\end{equation}

\noindent\textit{한 줄씩 읽기.}
\begin{itemize}[leftmargin=1.6em]
\item $\sum_i$ : 모든 삼각형 면에 대해 더한다.
\item $A_i$ : 그 면의 넓이로 \textbf{가중치}\footnote{가중치(weight): ``얼마나 비중
있게 칠지''를 정하는 곱하는 수. 큰 면일수록 전체 경향에 더 크게 반영되도록 면적 $A_i$를
곱해 줍니다. 학교 성적에서 중간고사보다 기말고사 비중을 크게 두는 것과 같은 원리입니다.}를
준다. 큰 면일수록 더 크게 반영.
\item $m_i m_i^{T}$ : 법선벡터 자기 자신을 ``표 형태''로 곱한 것
(\textbf{외적}\footnote{$m m^{T}$ (외적, outer product): 세로 벡터 $m$ 과 가로 벡터
$m^{T}$($T$는 \emph{전치}, 즉 세로를 가로로 눕힘\,\footnotemark)를 곱해
\emph{$3\times3$ 정사각형 표(행렬)}를 만드는 연산. 이 표 하나가 ``그 면이 가리키는
방향''의 정보를 담습니다. 내적($m\cdot m=$ 숫자 하나)과 헷갈리지 마세요. 외적은
\emph{표}, 내적은 \emph{숫자}입니다.}\footnotetext{전치(transpose, $T$): 행렬의 행과
열을 맞바꾸는 것. 세로로 선 숫자들을 가로로 눕히는 조작이라고 생각하면 됩니다.}).
이 작은 표들을 모두 더하면, 표면 법선이 \emph{주로 어느 방향들로 퍼져 있는지}를
요약하는 하나의 큰 표 $M$ 이 됩니다.
\end{itemize}

\begin{notebox}
\textbf{이 $M$ 이 바로 PCA와 같은 원리입니다.}
PCA\footnote{PCA(Principal Component Analysis, 주성분 분석): 여러 방향으로 흩어진
데이터에서 \emph{가장 많이 퍼진 주된 방향}을 찾아내는 통계 기법. 예를 들어 키-몸무게
점들이 비스듬한 타원으로 퍼져 있으면, PCA는 그 타원의 \emph{긴 축} 방향을 ``첫 번째
주성분''으로 찾아 줍니다. ``데이터가 가장 길게 늘어선 방향''을 찾는 도구라고 보면
됩니다.}는 데이터가 ``가장 많이 퍼진 방향''을 찾는 기법인데, 그 계산의 핵심이 바로
식 \eqref{eq:M}처럼 \emph{(데이터)$\times$(데이터)$^{T}$ 를 모두 더한 표}
(공분산행렬\footnote{공분산행렬(covariance matrix): 여러 변수(또는 방향)가 서로
얼마나 함께 변하는지를 한 표에 모은 것. PCA는 이 표의 \emph{고유벡터}를 구해 주성분
방향을 얻습니다.})를 만든 뒤, 그 표의 \textbf{고유벡터}\footnote{고유벡터
(eigenvector)\,·\,고유값(eigenvalue): 어떤 행렬 $M$ 을 곱해도 \emph{방향이 바뀌지
않고 길이만 변하는} 특별한 벡터가 고유벡터, 그때 늘어난 배율이 고유값입니다.
($M v = \lambda v$, 여기서 $v$=고유벡터, $\lambda$=고유값.) PCA에서 \emph{가장 큰
고유값}에 해당하는 고유벡터가 ``데이터가 가장 많이 퍼진 주된 방향''입니다.}를 구하는
것이기 때문입니다. 그래서 $M$ 의 고유벡터는 ``이 물체의 표면이 전체적으로 향하는
대표 방향''을 알려 줍니다.
\end{notebox}

\textbf{그런데 이걸로는 부족합니다.} $M$ 은 면들이 \emph{전체적으로} 어디를 보는지만
요약할 뿐, 정작 우리가 알고 싶은 것 ―
(1) 어떤 면이 \emph{특정 출력 방향에서} 오버행인가,
(2) 그 면을 \emph{실제로 받쳐 줄} 다른 면이 있는가,
(3) 못 받치면 바닥판까지 \emph{얼마나 먼가} ―
는 전혀 담지 못합니다. 게다가 $M$ 은 \textbf{방향을 바꿔 넣어도 답이 같은}
대칭적\footnote{대칭(symmetric): 여기서는 ``위/아래, 앞/뒤를 뒤집어도 결과가 똑같다''는
뜻. 하지만 3D 프린팅에서 \emph{위로} 향한 면과 \emph{아래로} 처진 면은 완전히 다릅니다
(아래로 처진 면만 서포트가 필요!). 그래서 위아래를 구별하지 못하는 대칭적인 $M$ 으로는
부족합니다.} 양이라, ``위로 향한 면''과 ``아래로 처진 면''을 구별하지 못합니다.
바로 이 약점을 고치려고 SFTF가 등장합니다.

%======================================================================
\section{SFTF의 핵심 수식 (한 줄씩)}
%======================================================================

이제 SFTF가 하나의 출력 방향 $n$ 에 대해 ``서포트가 얼마나 들지''를 어떻게
점수로 매기는지 차례로 봅니다. 핵심 발상은 \textbf{지지(받침)를 ``흐름''으로 보는}
것입니다 ― 마치 물이 위에서 아래로 흐르듯, ``받쳐야 할 무게''가 처진 면에서
아래의 받침면(또는 바닥)으로 흘러내린다고 상상합니다.

\subsection{1단계 — 이 면이 ``처진 면''인가? (오버행 계수)}

\begin{equation}
O_i(n)=\max\!\big(0,\;-m_i\cdot n\big)
\label{eq:overhang}
\end{equation}

\noindent\textit{한 줄씩 읽기.}
\begin{itemize}[leftmargin=1.6em]
\item $n$ : 우리가 시험 중인 \textbf{출력 방향}(위쪽 방향). $m_i$ : 면 $i$ 의 법선
(그 면이 바라보는 쪽).
\item $m_i\cdot n$ : 면이 \emph{출력 방향과 같은 쪽}을 보면 양수(+), \emph{반대쪽
(즉 아래로)}을 보면 음수($-$). (앞서 배운 ``내적 = 방향 비슷한 정도''!)
\item $-m_i\cdot n$ : 부호를 뒤집었으니, \textbf{아래로 처진 면일수록 큰 양수}가
됩니다.
\item $\max(0,\;\cdot\,)$ : 음수면 그냥 $0$ 으로 만든다는 뜻\footnote{$\max(0,x)$:
$x$ 가 양수면 그대로 두고, 음수면 0으로 바꾸는 함수. ``나쁜(음수) 값은 무시하고
좋은(양수) 값만 센다''는 장치. 머신러닝에서 ReLU라고도 부릅니다.}. 즉 \emph{위로
향한 면(서포트가 필요 없는 면)은 0점} 처리하고, \emph{아래로 처진 면만 점수}를 줍니다.
\end{itemize}

\begin{readbox}
\textbf{한 줄 요약.} $O_i(n)$ 은 ``면 $i$ 가 이 출력 방향에서 얼마나 심하게
아래로 처졌나''를 나타내는 점수입니다. $0$ 이면 안 처진 것, 클수록 심하게 처진
오버행입니다.
\end{readbox}

\subsection{2단계 — 누가 이 면을 받쳐 주나? (레이캐스트)}

처진 면 $i$ 를 발견했으면, ``그 아래에 받쳐 줄 다른 면이 있나?''를 봐야 합니다.
이를 위해 면 $i$ 에서 \emph{아래쪽}($-n$ 방향)으로 가상의 직선(광선)을 쏩니다. 이를
\textbf{레이캐스트}\footnote{레이캐스트(ray casting): 한 점에서 특정 방향으로
\emph{가상의 빛줄기(ray)}를 쏘아, 그 빛이 \emph{처음 부딪히는} 물체를 찾는 기법.
게임에서 ``총을 쏘면 무엇에 맞나'', 햇빛 그림자 계산 등에 쓰입니다. 여기서는
``처진 면에서 아래로 빛을 쏘면 어느 면이 받아 주나''를 찾습니다.}라고 합니다.
빛줄기가 처음 맞은 면을 받침면 $j$ 라고 합시다. 단, 아무 면이나 인정하지 않고
세 가지 조건을 통과해야 합니다.

\begin{itemize}[leftmargin=1.6em]
\item 자기 자신은 안 됨 ($j \neq i$).
\item 받침면이 처진 면보다 \textbf{아래에} 있어야 함 (높이차 $h_{ij}>0$).
\item 받침면이 \textbf{위를 보고 있어야} 함 ($m_j\cdot n>0.05$, 즉 어느 정도
하늘을 향해야 받칠 수 있음).
\end{itemize}

받침면을 찾았으면, \textbf{둘 사이의 높이차}와 \textbf{영향력}을 계산합니다.

\begin{equation}
h_{ij}(n)=(c_i-c_j)\cdot n
\label{eq:height}
\end{equation}
\noindent\textit{읽기:} $c_i, c_j$ 는 두 면의 \textbf{중심점}\footnote{중심점
(centroid): 삼각형의 한가운데 점(세 꼭짓점의 평균 위치). 면 전체를 대표하는 위치로
씁니다.} 위치입니다. 두 점의 위치 차이를 출력 방향 $n$ 에 내적했으니, $h_{ij}$ 는
``처진 면이 받침면보다 \emph{출력 방향으로 얼마나 높이} 떠 있나'' = 둘 사이의
세로 거리입니다.

\begin{equation}
w_{ij}(n)=\frac{O_i(n)}{1+\alpha\,h_{ij}(n)},
\qquad \alpha=\frac{1}{D}
\label{eq:w}
\end{equation}
\noindent\textit{한 줄씩 읽기.}
\begin{itemize}[leftmargin=1.6em]
\item 분자 $O_i(n)$ : 면이 심하게 처질수록(=받쳐야 할 게 많을수록) 영향력 $\uparrow$.
\item 분모 $1+\alpha h_{ij}$ : 받침면이 \emph{멀수록}(높이차 $h_{ij}$ 가 클수록)
분모가 커져서 영향력 $\downarrow$. 즉 \textbf{멀리 떨어진 받침은 덜 쳐준다}는 뜻.
가까이서 받쳐 주는 게 더 확실하니까요.
\item $\alpha=1/D$ : 여기서 $D$ 는 물체를 감싸는 상자의 \textbf{대각선
길이}\footnote{bbox 대각선(bounding-box diagonal): 물체를 딱 감싸는 직육면체 상자의
\emph{대각선} 길이. 물체의 ``전체 크기''를 나타내는 대표 길이입니다. 높이차 $h$ 를
이 크기로 나눠 줌으로써, 큰 물체든 작은 물체든 \emph{같은 기준}으로 거리를 비교할 수
있게 됩니다(단위 맞추기, 정규화\,\footnotemark).}\footnotetext{정규화
(normalization): 서로 크기가 다른 값들을 \emph{같은 잣대}로 비교할 수 있게 비율로
바꾸는 것. 예: 시험마다 만점이 다르면 100점 만점으로 환산해 비교.}. $h$ 를 $D$ 로
나눠 ``물체 크기 대비 얼마나 먼가''로 바꿉니다.
\end{itemize}

\begin{readbox}
\textbf{한 줄 요약.} $w_{ij}$ 는 ``처진 면 $i$ 가 받침면 $j$ 로부터 받는 받침의
세기''입니다. \emph{많이 처질수록} 크고, \emph{받침이 멀수록} 작아집니다.
\end{readbox}

\subsection{3단계 — 모든 받침 관계를 하나로 모으기 (흐름 텐서)}

이제 물체 안의 모든 (처진 면 $i$, 받침면 $j$) 짝에 대해 위 영향력을 모읍니다.

\begin{equation}
F(n)=\sum_{(i,j)\in\mathcal{R}(n)} w_{ij}(n)\,A_iA_j\,m_im_j^{T}
\label{eq:F}
\end{equation}
\noindent\textit{한 줄씩 읽기.}
\begin{itemize}[leftmargin=1.6em]
\item $\sum_{(i,j)\in\mathcal{R}(n)}$ : 레이캐스트로 찾은 \textbf{모든 유효한
받침 짝}\footnote{$\mathcal{R}(n)$: 출력 방향 $n$ 에서 ``처진 면 $\to$ 받침면''으로
인정된 (i,\,j) 짝들의 모임. 위 세 조건(자기 자신 아님, 아래에 있음, 위를 봄)을 통과한
쌍만 들어갑니다.}에 대해 더한다.
\item $w_{ij}$ : 방금 구한 받침 세기.
\item $A_iA_j$ : 두 면이 \emph{둘 다 클수록} 더 중요하게(면적 가중).
\item $m_im_j^{T}$ : 식 \eqref{eq:M}에서처럼 두 면의 방향 정보를 ``표''로 묶은 것.
단, 여기서는 \emph{처진 면}과 \emph{받침면}이 \textbf{서로 다른} 면이라
$m_i m_j^{T}$ 가 됩니다(자기 자신끼리인 $m_i m_i^T$ 가 아님).
\end{itemize}

\begin{notebox}
\textbf{``비대칭''이 핵심.} 앞의 $M$(식 \eqref{eq:M})은 위아래를 구별 못 하는
대칭이었습니다. 하지만 $F(n)$ 은 ``처진 면 $\to$ 받침면''이라는 \emph{방향}이
들어 있어서 \textbf{비대칭}\footnote{비대칭(asymmetric): ``A에서 B로''와 ``B에서
A로''가 다른 성질. 물이 높은 데서 낮은 데로만 흐르듯, 받침도 위에서 아래로만
일어납니다. 이 방향성을 담았기에 위/아래를 구별할 수 있습니다.}입니다. 그래서
출력 방향 $n$ 을 바꾸면 $F(n)$ 도 제대로 바뀝니다. 이름에 ``\textbf{흐름}(Flow)''이
붙은 이유입니다 ― 받침이 \emph{한 방향으로 흐르는} 구조를 담았으니까요.
\end{notebox}

\subsection{보너스 — 계산을 너무 무겁지 않게 (레이 예산)}

면이 100만 개인 큰 모델에서 모든 처진 면마다 빛을 쏘면 너무 느립니다. 그래서
\emph{쏘는 빛의 개수에 상한}을 둡니다.

\begin{equation}
K_{ray}=\mathrm{clip}\!\left(0.02\,N_f,\;500,\;3000\right)
\label{eq:kray}
\end{equation}
\noindent\textit{읽기:} $N_f$ 는 면의 개수. ``전체 면의 2\%(=$0.02 N_f$)만큼만 빛을
쏘되, 너무 적지도(최소 500개) 너무 많지도(최대 3000개) 않게 자른다''는 규칙입니다.
$\mathrm{clip}(x,a,b)$\footnote{clip(자르기): 값 $x$ 가 하한 $a$ 보다 작으면 $a$ 로,
상한 $b$ 보다 크면 $b$ 로 맞춰 \emph{$[a,b]$ 범위 안으로 가두는} 함수. 볼륨을
0$\sim$100 사이로만 조절되게 막아 두는 것과 같습니다.}이 그 ``범위 안으로 가두기''
함수입니다. 덕분에 모델이 아무리 커져도 한 방향당 계산량이 \emph{3000개를 넘지 않아}
큰 모델도 감당할 수 있습니다.

%======================================================================
\section{가장 중요한 항: 바닥판(Ground Node)}
%======================================================================

처진 면이 \emph{아래로 빛을 쐈는데 받쳐 줄 면이 하나도 없으면} 어떻게 될까요?
결국 \textbf{프린터 바닥판}(build plate)이 맨 밑에서 받쳐 줘야 합니다. 이 ``바닥에서
올라오는 긴 기둥''이 사실 서포트의 가장 큰 부분이라, 이 항이 \textbf{가장 중요}합니다.

\subsection{바닥판 받침의 양}

\begin{equation}
z_{plate}(n)=\min_{v\in V} v\cdot n,
\qquad
h_{iP}(n)=c_i\cdot n - z_{plate}(n)
\label{eq:plate-height}
\end{equation}
\noindent\textit{한 줄씩 읽기.}
\begin{itemize}[leftmargin=1.6em]
\item $z_{plate}=\min_{v\in V} v\cdot n$ : 모든 꼭짓점\footnote{$V$ 는 메쉬의 모든
꼭짓점(vertex)들의 모임, $v$ 는 그중 하나. $\min$(최소)은 ``가장 작은 값''을 고르라는
뜻.} 중 출력 방향으로 \textbf{가장 낮은 위치} = 바닥판이 닿는 높이.
\item $h_{iP}=c_i\cdot n - z_{plate}$ : 처진 면 $i$ 가 \textbf{바닥에서 얼마나 높이}
떠 있나 = 바닥까지 세워야 할 기둥의 길이.
\end{itemize}

\begin{equation}
B(n)=\sum_{i\in\mathcal{B}(n)} O_i(n)\,A_i\,\big(1+\alpha\,h_{iP}(n)\big)
\label{eq:B}
\end{equation}
\noindent\textit{한 줄씩 읽기.}
\begin{itemize}[leftmargin=1.6em]
\item $\sum_{i\in\mathcal{B}(n)}$ : \textbf{받침면을 못 찾은} 처진 면들만 모아
더한다\footnote{$\mathcal{B}(n)$: 레이캐스트로 받침면을 찾지 못해 \emph{결국 바닥판이
받쳐야 하는} 처진 면들의 모임.}.
\item $O_i(n)\,A_i$ : 심하게 처질수록, 면이 클수록 비용 $\uparrow$.
\item $(1+\alpha h_{iP})$ : 식 \eqref{eq:w}의 분모와 닮았지만 여기서는 \emph{곱}으로
들어갑니다. \textbf{바닥에서 높이 떠 있을수록 더 긴 기둥}이 필요하니 비용이 커지는
게 맞습니다.
\end{itemize}

\begin{readbox}
\textbf{한 줄 요약.} $B(n)$ 은 ``바닥판에서 직접 받쳐 올려야 하는 서포트의 총량''
입니다. 뒤에서 보겠지만 \textbf{이 항이 서포트 양을 가장 잘 예측}합니다.
\end{readbox}

\subsection{멋진 통합: 바닥판도 ``가상의 받침면''으로 보기}

논문의 핵심 아이디어 중 하나는, 위에서 따로 만든 바닥판 항 $B(n)$ 이 사실
``\textbf{법선이 출력 방향 $n$ 인 가상의 받침면}을 하나 추가''한 것과 \emph{정확히
같다}는 점입니다. 즉 식 \eqref{eq:F}의 흐름 텐서에 바닥판 몫을 더해 줍니다.

\begin{equation}
\tilde F(n)=\underbrace{\sum_{(i,j)\in\mathcal{R}(n)} w_{ij}(n)A_iA_jm_im_j^{T}}_{\text{면끼리 받침(식 }\eqref{eq:F}\text{)}}
+\underbrace{\sum_{i\in\mathcal{B}(n)} A_i(1+\alpha h_{iP})\,m_i n^{T}}_{\text{바닥판 = 가상 받침면 }n}
\label{eq:Ftilde}
\end{equation}
\noindent\textit{읽기:} 두 번째 합에서 받침면의 방향 자리에 \emph{진짜 면 $m_j$ 대신
출력 방향 $n$} 이 들어간 것($m_i n^{T}$)만 다릅니다. ``바닥판은 출력 방향에서 봤을 때
완벽히 위를 보는 받침면''이라고 본 것입니다.

\subsection{레일리 값: 텐서를 하나의 비용 숫자로}

표(텐서) $\tilde F$ 를 다시 \emph{하나의 비용 점수}로 바꿔야 비교가 됩니다. 이때
쓰는 게 \textbf{레일리 값}\footnote{레일리 몫(Rayleigh quotient): 정사각형 표(행렬)
$F$ 와 방향 $n$ 에 대해 $n^{T} F n$ 으로 계산하는 \emph{하나의 숫자}. 직관적으로는
``그 표가 방향 $n$ 쪽으로 얼마나 강하게 작용하나''를 재는 값입니다. 식의 양쪽에서
$n$ 으로 ``눌러 짜낸다(contraction)''고 생각하면 됩니다. PCA의 고유값도 사실 이
레일리 값을 가장 크게/작게 만드는 방향에서 나옵니다 ― 그래서 \eqref{eq:M}의 PCA
이야기와 한 가족입니다.}입니다.

\begin{equation}
\tilde R(n)=\max\!\big(0,\;-\,n^{T}\tilde F_s(n)\,n\big)=R(n)+B(n)
\label{eq:Rtilde}
\end{equation}
\noindent\textit{한 줄씩 읽기.}
\begin{itemize}[leftmargin=1.6em]
\item $n^{T}\tilde F_s\, n$ : 통합 텐서를 출력 방향 $n$ 으로 양쪽에서 눌러 짠
\emph{하나의 숫자}\footnote{$\tilde F_s$ 의 아래첨자 $s$ 는 ``대칭 부분(symmetric
part)''만 쓴다는 표시입니다. 레일리 값 계산에는 표의 대칭 부분만 기여하기 때문입니다.
세부는 넘어가도 좋습니다.}.
\item 앞의 마이너스와 $\max(0,\cdot)$ : 식 \eqref{eq:overhang}과 같은 정신 ―
``서포트가 필요한(나쁜) 양만 0 이상으로 센다''.
\item \textbf{$=R(n)+B(n)$} : 놀랍게도, 이렇게 통합 텐서로 한 번에 계산한 값이
``면끼리 받침 비용 $R$'' 더하기 ``바닥판 비용 $B$''와 \emph{정확히} 같습니다.
\end{itemize}

\begin{notebox}
\textbf{왜 이게 중요한가.} $B$ 를 그냥 ``땜질로 끼워 넣은 보정값''이 아니라,
\emph{같은 텐서 수학에서 자연스럽게 나오는 한 부분}임을 보인 것입니다. 논문은 이
등식이 컴퓨터 계산에서도 소수점 아래 $10^{-12}$(0.000000000001) 수준까지 \emph{정확히}
맞아떨어짐을 확인했습니다. 이것이 이 방법을 ``\textbf{텐서}장''이라 부를 자격을 줍니다.
\end{notebox}

%======================================================================
\section{최종 점수와 후보 고르기}
%======================================================================

\subsection{기본 점수 \texorpdfstring{$J$}{J}}

실제로 쓰는 기본 점수는 다음과 같습니다.

\begin{equation}
J(n)=\tilde R(n)+P(n),\qquad
P(n)=\sum_{(i,j)\in\mathcal{R}(n)}O_i(n)\,A_iA_j
\label{eq:J}
\end{equation}
\noindent\textit{읽기:} $\tilde R=R+B$ 는 방금 구한 ``레일리 받침 비용(면끼리 + 바닥판)''.
$P$ 는 ``면끼리 받침의 양''을 따로 센 항입니다(높이 감쇠를 뺀 순수한 오버행$\times$면적의
합). 둘을 더해 ``이 방향에서 서포트가 대략 얼마나 들지''를 한 숫자 $J(n)$ 으로 냅니다.
$J$ 가 \emph{작을수록} 좋은(서포트 적은) 방향입니다.

\subsection{후보를 만드는 절차}

\begin{enumerate}[leftmargin=1.8em]
\item \textbf{넓게 훑기:} 구 표면에 고르게 뿌린 512개 방향
(\textbf{피보나치 구}\footnote{피보나치 구(Fibonacci sphere): 공 표면에 점들을
\emph{최대한 고르게} 흩뿌리는 잘 알려진 방법. 해바라기 씨앗이 촘촘하고 균일하게 박힌
배열과 같은 수학(황금비)을 씁니다. 덕분에 적은 점으로도 모든 방향을 골고루 시험할 수
있습니다.})에서 $J$ 를 계산.
\item \textbf{좁혀 들어가기:} 점수가 좋았던 방향들 주변 $2^\circ\sim8^\circ$ 의
좁은 원뿔 안을 더 촘촘히 다시 계산.
\item \textbf{겹치는 후보 정리:} 거의 같은 방향이 여러 개 뽑히면 하나만 남기는
\textbf{비최대 억제}\footnote{비최대 억제(NMS, Non-Maximum Suppression): 비슷한 위치에
여러 후보가 몰려 있을 때 \emph{가장 좋은 것 하나만} 남기고 나머지는 지우는 정리 기법.
``서로 충분히 떨어진'' 다양한 후보들만 남겨, 한곳에 쏠리지 않게 합니다.}로 서로
떨어진 대표 후보들만 남김.
\end{enumerate}

\subsection{순위 정규화로 최종 정렬 (단위가 다른 값들을 공정하게)}

마지막에는 $[\,J,R,P,B,H,S,\sigma_1\,]$ 같은 여러 특징\footnote{여기서 $H$=받침을
찾은 횟수, $S$\,·\,$\sigma_1$=텐서에서 나오는 보조 수치(특잇값\,\footnotemark)입니다.
세부는 몰라도 됩니다 ― 핵심은 ``여러 종류의 단서를 함께 본다''는 점입니다.}\footnotetext{특잇값
(singular value): 행렬이 방향마다 얼마나 늘이고 줄이는지를 나타내는 수치들. 가장 큰
것을 $\sigma_1$ 로 적습니다.}을 함께 써서 후보 순위를 매깁니다. 그런데 이 값들은
\emph{단위가 제각각}(어떤 건 부피, 어떤 건 개수)이라 그냥 더하면 불공정합니다.
그래서 각 특징을 \textbf{순위(등수)로 바꿔서} 비교합니다.

\begin{notebox}
\textbf{순위 정규화(rank normalization)란?} 키 175\,cm, 몸무게 70\,kg처럼 단위가
다른 값을 직접 더할 순 없습니다. 대신 ``키 등수 3등, 몸무게 등수 5등''처럼
\emph{등수}로 바꾸면 단위가 사라져 공정하게 합칠 수 있습니다\footnote{이렇게 값을
\emph{등수}로 바꿔 분석하는 것은 통계학의 표준 기법입니다. 뒤에 나오는 ``스피어만
상관''도 같은 정신입니다.}. SFTF도 각 특징을 후보들 사이의 등수로 바꾼 뒤 가중합해
최종 순위를 정합니다. (이는 정밀한 기하 최적화가 아니라, 경험적으로 좋은 순서를
찾는 \emph{보정}입니다.)
\end{notebox}

%======================================================================
\section{후보를 ``검증''하고, 어려우면 ``더 본다'' (AVE)}
%======================================================================

\subsection{국소 정밀검증}

SFTF가 추천한 상위 $K$ 개 방향 각각의 \textbf{둘레 $\pm10^\circ$ 만} 정밀 계산기
(TOMO\footnote{TOMO: 이 연구에서 \emph{진짜 서포트 양}을 정밀하게 재는 기준 계산기
(검증용\,·\,비교용). CPU판과 GPU(CUDA)판이 있습니다. SFTF는 이 느린 계산기에 넘길
방향의 \emph{개수를 줄여 주는} 역할입니다.})로 꼼꼼히 검사하고, 그중 가장 좋은 값을
고릅니다. ``전부 다''가 아니라 ``유망한 곳 둘레만'' 보는 게 빠름의 비결입니다.

\begin{readbox}
\textbf{핵심 성과.} 상위 \textbf{20개} 후보의 둘레만 검증해도, 전체 격자의
\textbf{단 3.57\%}만 보고서 최적 대비 평균 \textbf{1.13배}(최악 1.58배)의 서포트
양을 찾아냈습니다. 같은 예산을 ``아무 데나/규칙적으로'' 쓴 비교군(평균 1.64\,$\sim$\,2.62배)보다
훨씬 좋습니다 ― \emph{어디를 볼지 잘 찍는다}는 증거입니다.
\end{readbox}

\subsection{어려운 경우 자동 확대 (AVE)}

가끔 정답이 아주 좁은 골짜기에 숨어 있어, 상위 10개만 봐서는 놓치는 ``어려운
케이스''가 있습니다(예: \emph{Happy} 모델). 이를 위해 \textbf{AVE}\footnote{AVE
(Adaptive Verification Escalation, 적응적 검증 확대): ``검사해 보니 확신이 안 서면,
\emph{그때만} 검사 범위를 자동으로 넓히는'' 규칙. 평소엔 적게 보고, 미심쩍을 때만 더
보는 똑똑한 절약 장치입니다.}라는 규칙을 둡니다. 검증 결과가 미덥지 않다는 신호가
잡힐 때만 자동으로 후보를 더 넓혀 검사합니다.

\begin{equation}
\begin{aligned}
\mathrm{trigger}_{\mathrm{AVE}}
&=\big([\eta\ge0.03]\wedge[\mathrm{edge}=1\vee g\le1.5]\big)\;\vee\;[\mathrm{edge}=1\wedge g\ge3.0],\\[2pt]
\eta&=\frac{V_{ss}^{\mathrm{loc}}}{V_{\mathrm{bbox}}},\qquad
g=\frac{V_{ss}^{\mathrm{seed}}}{V_{ss}^{\mathrm{loc}}}
\end{aligned}
\label{eq:ave}
\end{equation}
\noindent\textit{한 줄씩 읽기} (기호 $\wedge$=``그리고(and)'', $\vee$=``또는(or)''\footnote{논리
기호: $\wedge$ 는 두 조건이 \emph{모두} 참일 때 참(``그리고''), $\vee$ 는 \emph{하나라도}
참이면 참(``또는''). $[\,\cdot\,]$ 는 ``안의 조건이 맞으면 1(참), 아니면 0(거짓)''.}):
\begin{itemize}[leftmargin=1.6em]
\item $\eta=V_{ss}^{\mathrm{loc}}/V_{\mathrm{bbox}}$ : 찾은 서포트 양을 물체 상자
부피로 나눈 ``상대적 크기''. $\eta\ge0.03$ 이면 ``서포트가 꽤 많네?'' 하는 신호.
\item $g=V_{ss}^{\mathrm{seed}}/V_{ss}^{\mathrm{loc}}$ : 처음 찍은 값 대비 검증 후
값의 비율. 검증으로 많이 좋아졌다면($g$ 큼) 처음 추천이 부실했다는 뜻.
\item $\mathrm{edge}=1$ : 찾은 최적 칸이 검사 \emph{창문 가장자리}에 걸쳐 있음 =
``창문 밖에 더 좋은 게 있을지도?'' 하는 의심 신호.
\item 전체: 이런 의심 신호들이 조합되면 \textbf{검사 범위를 자동 확대}.
\end{itemize}

\begin{notebox}
\textbf{AVE는 손해 볼 일이 없습니다.} 확대는 \emph{후보를 더 볼 뿐}이라, 예산
(계산량)은 더 쓸 수 있어도 \emph{이미 찾은 답이 나빠질 수는 없습니다}. 실제로 어려운
케이스의 최악 비율을 \textbf{8.52배 $\to$ 1.06배}로 끌어내렸고(평균 예산은 3\,$\sim$\,4.5\%),
다섯 그룹 점검에서도 최악 \textbf{17.57배 $\to$ 1.37배}로 좋아졌습니다.
\end{notebox}

%======================================================================
\section{결과 요약과 정직한 한계}
%======================================================================

\subsection{얼마나 빠른가}

같은 1$^\circ$ 격자 조건에서, SFTF는 정밀 계산기(TOMO\_CPU) 대비
\textbf{파이썬 구현으로 8\,$\sim$\,382배}, \textbf{C++ 구현으로 최대 5만 6천 배}까지
빨랐습니다. 빠름의 핵심은 ``계산기를 바꿔서''가 아니라 ``\emph{검사할 방향의 개수
자체를 확 줄여서}''입니다. 식 \eqref{eq:kray}의 레이 예산 덕분에 면이 100만 개인
거대 모델도 감당했습니다.

\subsection{어떤 항이 가장 쓸모 있나}

여러 특징과 ``진짜 서포트 양''의 \textbf{스피어만 상관}\footnote{스피어만 상관계수
(Spearman rank correlation): 두 양이 ``\emph{같이 커지고 같이 작아지는지}''를
\emph{등수 기준}으로 재는 값($-1\sim1$). $+1$ 에 가까우면 ``한쪽이 크면 다른 쪽도
크다''(강한 양의 관계), $0$ 이면 무관, $-1$ 이면 정반대. 앞서 본 순위 정규화처럼
값을 등수로 바꿔 보기 때문에, 단위나 곡선 모양에 휘둘리지 않고 ``경향''만 봅니다.}을
재 보니, \textbf{바닥판 항 $B$ 가 가장 강한 단일 예측자}였습니다(평균 0.58). 통합
레일리 항 $\tilde R=R+B$ 는 더 높은 0.60이었고요. 이는 ``\emph{바닥판 받침이 서포트
양을 지배한다}''는 직관과 정확히 맞습니다.

\subsection{정직한 한계 (이게 제일 중요)}

\begin{notebox}
\textbf{SFTF는 ``정밀 최적화기''가 아닙니다.} 순위만으로 정답을 내는 점수는, 한 번도
본 적 없는 새 모델에서는 가끔 크게 빗나갑니다(교차검증\footnote{교차검증
(cross-validation): 가진 데이터의 일부를 \emph{일부러 가려 두고} 나머지로만 규칙을
만든 뒤, 가려 둔 것으로 시험해 ``처음 보는 데이터에도 통하는가''를 점검하는 방법.
시험 범위만 외운 게 아니라 진짜 실력인지 확인하는 모의고사와 같습니다. 특히 한 모델씩
빼고 시험하는 방식을 LOMO\,\footnotemark 라고 합니다.}\footnotetext{LOMO
(Leave-One-Mesh-Out): 메쉬(모델)를 하나씩 차례로 빼 두고 나머지로 학습\,·\,그 하나로
시험하기를 반복하는 교차검증 방식.}에서 최악 비율이 100배를 넘는 경우도 있었습니다).
\end{notebox}

그래서 이 방법의 올바른 사용법은 처음 말한 그대로입니다 ― \textbf{값싼 후보 추천
(SFTF) $\to$ 유망한 곳만 정밀 검증(TOMO) $\to$ 어려우면 자동 확대(AVE)}. 이 조합으로만
신뢰할 수 있는 결과가 나옵니다. SFTF는 정밀 설계를 \emph{대체}하는 게 아니라, 비싼
검증을 \emph{어디에 집중할지} 빠르게 알려 주는 똑똑한 \textbf{출발점(warm start)}입니다.

\begin{readbox}
\textbf{전체 한 줄 정리.} ``서포트 적은 출력 방향 찾기''라는 비싼 문제에서, SFTF는
받침을 \emph{흐름}으로 보는 텐서로 유망한 방향을 \textbf{빠르고 값싸게} 추려 주고
(특히 바닥판 항 $B$ 가 핵심), 그 둘레만 정밀 검증하며, 어려운 경우엔 자동으로 더
들여다봅니다. 똑똑하게 ``어디를 볼지'' 잘 찍어 주는 도구입니다.
\end{readbox}

\vfill
\begin{center}\small
\textit{원문: \texttt{SFTF\_draft.tex} (Support Flow Tensor Field for Fast
Build-Orientation Candidate Generation).\\
이 글은 같은 내용을 고등학생 눈높이로 풀어 쓴 해설본입니다.}
\end{center}

\end{document}
